% ======================================================================
% 1. RESUMEN EJECUTIVO
% ======================================================================
\begin{abstractd}
	Este reporte documenta el proceso de investigación y planificación para la Tarea 3 del curso Desafíos de Innovación en Ingeniería y Ciencias. El proyecto aborda la problemática de la contaminación por macroplásticos enterrados en las playas chilenas, un residuo que escapa a los métodos de limpieza manual tradicionales. A través de una revisión bibliográfica exhaustiva, se determinó que la solución óptima requiere de un sistema de cribado mecánico. Se propone el desarrollo de "Wallis", un prototipo robótico autónomo diseñado para operar de forma eficiente y silenciosa. El documento detalla la metodología de búsqueda de información, el análisis ético de la solución, y la planificación de una observación de campo en la zona central de Chile para validar las hipótesis de diseño. Los resultados preliminares indican que el proyecto es técnicamente factible y socialmente necesario para mitigar el impacto ambiental y económico en las zonas costeras.
\end{abstractd}

\newpage
% ======================================================================
% 2. INTRODUCCIÓN
% ======================================================================
\section{Introducción}
La contaminación ambiental en el borde costero chileno representa uno de los desafíos más críticos para la ingeniería ambiental contemporánea. Chile, con más de 6.500 kilómetros de costa, enfrenta una acumulación persistente de residuos sólidos, principalmente macroplásticos (colillas, tapas, envases), que son enterrados por la dinámica de la arena y el viento.

La limpieza manual, método predominante en las municipalidades del país, resulta ineficiente para extraer basura subsuperficial, lo que genera riesgos de accidentes para turistas y la degradación irreversible de los ecosistemas marinos. Este proyecto propone la creación de un robot autónomo que realice tamizado profundo durante la noche, optimizando la gestión de limpieza sin interferir con las actividades recreativas de la zona.

\subsection{Objetivos del Proyecto}
\begin{itemize}
	\item \textbf{Objetivo General:} Diseñar y validar un prototipo de robot autónomo capaz de recolectar residuos enterrados en arena de playa mediante tamizado mecánico.
	\item \textbf{Objetivos Específicos:}
	      \begin{enumerate}
		      \item Realizar un catastro de tipos y profundidades de basura mediante observación de campo.
		      \item Identificar las necesidades de los residentes y turistas a través de herramientas de diseño empático.
		      \item Evaluar la factibilidad técnica y el impacto ético de la automatización en tareas de limpieza pública.
	      \end{enumerate}
\end{itemize}

\newpage
% ======================================================================
% 3. DESARROLLO
% ======================================================================
\section{Desarrollo}

\subsection{Metodología de Búsqueda de Información}
Para la fundamentación técnica del proyecto, se realizó una revisión del estado del arte utilizando bases de datos académicas (Google Scholar, ScienceDirect) y repositorios institucionales de la Armada de Chile (DIRECTEMAR) y el Ministerio del Medio Ambiente (MMA). Las palabras clave empleadas fueron "cribado mecánico de arena", "beach cleaning robots", y "microplásticos en costas chilenas". Se analizaron 15 fuentes técnicas para asegurar la viabilidad de los materiales (PETG y Acero 304) y el sistema de tracción por orugas.

\subsection{Definición de la Problemática y Sintaxis}
Basado en los hallazgos iniciales, se ha definido el problema central del proyecto utilizando la sintaxis requerida:

\begin{center}
	\fbox{
		\begin{minipage}{0.8\linewidth}
			\centering \vspace{5pt}
			\textbf{Residentes y turistas de playas chilenas} (Usuario) + \textbf{necesitan playas libres de basura enterrada} (Necesidad) + \textbf{debido al riesgo de accidentes y degradación ambiental} (Síntoma) + \textbf{puesto que la limpieza manual municipal no logra extraer macroplásticos subsuperficiales} (Hallazgo).
			\vspace{5pt}
		\end{minipage}
	}
\end{center}

\subsection{Revisión del Estado del Arte y Referentes}
La recolección nacional en 2024 sumó 86 toneladas de basura, donde el 61\% corresponde a plásticos. A nivel internacional, países como España y Brasil ya implementan maquinaria de tamizado adaptada a tractores. En el ámbito robótico, existen soluciones limitadas:
\begin{itemize}
	\item \textbf{BeBot:} Excelente filtrado pero requiere control manual constante.
	\item \textbf{Robot Zambrano:} Autónomo mediante IA, pero limitado a recolección superficial (pala).
	\item \textbf{VERO:} Robot cuadrúpedo eficiente, pero especializado exclusivamente en colillas.
\end{itemize}
Wallis busca integrar lo mejor de estos mundos: autonomía total y capacidad de tamizado mecánico profundo.

\subsection{Estudio Ético del Problema}
Desde la perspectiva de la ingeniería, es fundamental considerar el impacto social de la automatización. El robot Wallis no busca desplazar al personal de aseo municipal, sino servir como una herramienta de aumento de capacidades, encargándose de los residuos que el ojo humano no puede detectar. Éticamente, se prioriza el funcionamiento nocturno para respetar la privacidad de los usuarios y minimizar la contaminación acústica. Asimismo, el uso de materiales reciclables (PETG) y sistemas eléctricos reduce la huella de carbono de la operación.

\subsection{Análisis Empático y AEIOU}
El análisis de los usuarios (residente Ana Jer, 56 años) indica una alta frustración por la ineficiencia municipal. La Matriz AEIOU revela que el entorno se ve saturado por el comercio ambulante que genera envases plásticos sin puntos de disposición cercanos.
\insertimage{img/2.jpg}{width=0.4\linewidth}{Evidencia de basura acumulada en el borde costero analizado.}

\newpage
% ======================================================================
% 4. CONCLUSIONES
% ======================================================================
\section{Conclusiones y Próximos Pasos}
Como equipo, el principal aprendizaje ha sido la comprensión de que la ingeniería debe ser sistémica: no basta con un mecanismo de filtrado eficiente si no se considera el comportamiento humano y la falta de fiscalización de leyes como la 21.100.

\subsection{Conexión con los ODS}
El proyecto se alinea directamente con:
\begin{itemize}
	\item \textbf{ODS 14 (Vida Submarina):} Prevención de la entrada de macro y microplásticos al océano.
	\item \textbf{ODS 12 (Producción y Consumo Responsables):} Al fomentar una gestión eficiente de los residuos urbanos en playas.
\end{itemize}

\subsection{Factibilidad y Avances}
Se concluye que el proyecto es factible técnica y económicamente mediante el uso de hardware abierto (Arduino/Raspberry Pi). Los próximos pasos consisten en la ejecución de la observación de campo planificada y el testeo del mecanismo de cribado en arena seca y húmeda.

\newpage
% --- REFERENCIAS ---
\nocite{*}
\insertindextitlepage{\namereferences}
\bibliography{library}

\newpage
% --- ANEXOS ---
\begin{appendixd}
	\section{Planificación Detallada de Observación de Campo}
	Se realizará un recorrido de 200 metros lineales en la playa principal de El Quisco. Se medirá la profundidad de los residuos encontrados utilizando una sonda manual y se registrará la frecuencia de paso de los camiones de basura municipales.

	\section{Matriz AEIOU (Registro Fotográfico)}
	\insertimage{img/1.jpg}{width=0.5\linewidth}{Limpieza manual observada durante la fase piloto.}
	\insertimage{img/3.jpg}{width=0.5\linewidth}{Interacciones negligentes de turistas con el entorno.}
	\insertimage{img/4.jpg}{width=0.5\linewidth}{Tipos de objetos (latas y vidrios) predominantes.}
	\insertimage{img/5.jpg}{width=0.5\linewidth}{Zonas de alta concurrencia donde se aplicará Wallis.}

	\section{Transcripción de Entrevista}
	\textbf{Entrevistada:} Ana Jer.
	\begin{itemize}
		\item \textit{"La municipalidad pasa diario, pero no dan abasto, mientras más escarbas la arena, más sale basura antigua."}
	\end{itemize}
\end{appendixd}
